\documentclass{article}
\usepackage[UTF8]{ctex}
\usepackage{amsmath,amssymb,amsfonts}
\usepackage{enumitem}
\usepackage{geometry}
\geometry{a4paper, margin=1in}
% === 错题本排版优化（不参与题目标识） ===
\setlength{\parindent}{0pt}
\setlength{\parskip}{0.35em plus 0.10em minus 0.08em}
\setlength{\emergencystretch}{3em}
\setlist{itemsep=0.25em, topsep=0.40em, parsep=0pt}
\allowdisplaybreaks[2]
% === 错题本排版优化结束 ===
\begin{document}
\section*{北京交通大学2024年数学分析真题}
\section*{5．北京交通大学2024年数学分析试卷（回忆版）}
一．数列 \(\left\{a_{n}\right\}\) 没有最大值，没有最小值，证明：\(\left\{a_{n}\right\}\) 发散．

二．\(f(x)\) 在 \([0,2]\) 上连续、在 \((0,2)\) 上可导，且 \(f(0)+f(1)=2, f(2)=1\) 。证明：存在 \(\xi \in(0,2)\) 使得 \(f^{\prime}(\xi)=0\) 。

三．计算 \(\int \frac{x \mathrm{e}^{x}}{(x+1)^{2}} \mathrm{~d} x\) ．

\subsection{题目7}
计算 \(\int \frac{1}{(x-1)^{2}(x-2)} \mathrm{d} x\) ．

四．设 \(f(x)\) 在 \([a, b]\) 上三阶可导，证明存在 \(\xi \in(a, b)\) 使得

\[
f(b)=f(a)+\frac{1}{2}(b-a)\left[f^{\prime}(a)+f^{\prime}(b)\right]-\frac{1}{12}(b-a)^{3} f^{\prime \prime \prime}(\xi) .
\]

\subsection{题目1}
1．\(f(x)=\frac{|\sin x|}{x}\) 在 \((-1,0)\) 和 \((0,1)\) 上一致连续．

\subsection{题目2}
2．\(f(x)=\frac{|\sin x|}{x}\) 在 \((-1,0) \cup(0,1)\) 上不一致连续．

六．设

\[
f(x)= \begin{cases}x y \sin \frac{1}{x^{2}+y^{2}}, & x^{2}+y^{2} \neq 0 \\ 0, & x^{2}+y^{2}=0\end{cases}
\]

\subsection{题目3}
1．求 \(\frac{\partial f}{\partial x}, \frac{\partial f}{\partial y}\) ．\\

\subsection{题目4}
2．\(\frac{\partial f}{\partial x}, \frac{\partial f}{\partial y}\) 在 \((0,0)\) 是否连续，\(f(x, y)\) 在 \((0,0)\) 是否可微？

七．证明 \(\left\{\begin{array}{l}x+y+u+v=0 ; \\ x^{2}+y^{2}+u^{4}+v^{4}=1 .\end{array}\right.\) 存在 \(\binom{u}{v}=\binom{u(x, y)}{v(x, y)}\) ，并计算 \(\frac{\partial u}{\partial x}, \frac{\partial^{2} u}{\partial x^{2}}\) ．

九．求曲线积分问题，应该是考察斯托克斯公式的应用．

\begin{itemize}
  \item[十．] 设 \(f_{n}(x)=\frac{n x}{1+\left(1+n^{2}\right) x^{2}}\) ．

  \item[1．] 证明：\(\left\{f_{n}(x)\right\}\) 在 \((0,1]\) 上不一致收敛。

  \item[2．] 证明：\(\left\{f_{n}(x)\right\}\) 在 \([\alpha, 1]\) 上一致收敛，其中 \(0<\alpha<1\) 。
\end{itemize}

十一．计算 \(\sum_{n=1}^{\infty} \frac{1}{(n+1) 2^{n}}\) ．
\end{document}
